\documentclass[11pt,a4paper]{article}
\usepackage[margin=27mm]{geometry}
\usepackage{amsmath,amssymb,amsthm,hyperref,longtable,booktabs,array}
\usepackage[T1]{fontenc}
\usepackage{lmodern}
\emergencystretch=1em
\providecommand{\tightlist}{\setlength{\itemsep}{0pt}\setlength{\parskip}{0pt}}
\newcommand{\pandocbounded}[1]{#1}
\title{Subfield and torus geometries}
\author{Human Compression project}\date{}
\begin{document}\maketitle
We realize every admissible parameter tuple in Constructions 8 and 9 of
the fixed Suzuki source. A valuation calculation gives all
intersections, including the cases in which a selected prime divides the
base torus order. Descent and cyclic coset coordinates prove flag
transitivity; the same intersections give thickness, residual
connectedness and residual weak primitivity.

\section{Source parameters and the base
field}\label{source-parameters-and-the-base-field}

Write the field exponent as \(n\), so the ambient group is
\(G=\operatorname{Sz}(2^n)\). The source chooses \(k=\alpha-1\) pairwise
nonisomorphic maximal simple subfield groups
\(\operatorname{Sz}(2^{n/a_i})\), where the relative degrees \(a_i\) are
distinct odd primes. Set \[
e=n/\prod_i a_i,\qquad t=2^e.
\]

The upper bound \(\alpha\le\sum_j e_j=\Omega(n)\) in both constructions
implies \(\Omega(e)=\Omega(n)-k\ge1\). Since \(n\) is odd,
\textbf{\(e\ge3\) and \(t\ge8\)}. Thus neither construction includes a
nonsimple base \(\operatorname{Sz}(2)\). The notation \(\Omega\) here
counts prime factors with multiplicity. Some \(a_i\) may also divide
\(e\); the proof allows that.

Construction 8 chooses distinct primes \[
\ell_i\mid\frac{t^{a_i}-1}{t-1}
\] and adds the dihedral subgroup of order \[
2(t-1)\prod_i\ell_i.
\]

Construction 9 chooses a base nonsplit torus order
\(\gamma=t+\sqrt{2t}+1\) or \(t-\sqrt{2t}+1\), and distinct primes
\(\ell_i\) such that \(\gamma\ell_i\) divides an appropriate nonsplit
torus order of \(\operatorname{Sz}(t^{a_i})\). It also requires
\(\gamma\prod_i\ell_i\) to divide one of the two ambient nonsplit torus
orders. Its added subgroup has order \[
4\gamma\prod_i\ell_i
\] and is cyclic by \(C_4\).

The source does \textbf{not} require \(\ell_i\nmid\gamma\). This note
covers those cases in which a selected prime divides the base order
explicitly. The arbitrary-rank existence theorem in the companion
descent paper uses the extra restriction \(a_i\nmid\tau_e\) only to
prove an unconditional supply of parameters; that restriction is not
needed for the present realization theorem for every admissible
parameter tuple.

\section{Both nonsplit branches have compatible binary
models}\label{both-nonsplit-branches-have-compatible-binary-models}

The companion descent paper constructs a torus \(\mathbf T\), stable
under the exceptional endomorphism \(\rho\), with \[
|\mathbf T^{\rho^d}|=\tau_d
=2^d+1-\big((1+i)^d+(1-i)^d\big)
\] for every odd \(d\). To obtain the other branch, retain its binary
order-four matrix \(r\) and antidiagonal \(w\), and set \[
c=r^2w.
\]

The retained exact seed check gives \[
|c|=5,\qquad \rho(c)=c,\qquad
C_{\operatorname{Sz}(2)}(c)=\langle c\rangle.
\] It also verifies that \(r\) conjugates \(c\) to \(c^2\) or \(c^3\),
so normalizes its centralizer. As the multiplicative order of 2 modulo 5
is four, \(c\) has four distinct eigenlines cyclically permuted by
binary Frobenius. Therefore \(\mathbf T'=C_{\operatorname{Sp}_4}(c)\) is
another maximal torus with the same four-cycle Frobenius calculation as
in the companion descent paper. It is \(\rho\)-stable because
\(\rho(c)=c\).

The character-lattice endomorphism has determinant two and trace
\(\pm2\); its five \(\rho\)-fixed points select trace \textbf{minus
two}. For odd \(d\), this gives \[
|(\mathbf T')^{\rho^d}|=\tau'_d
=2^d+1+\big((1+i)^d+(1-i)^d\big).
\] Both finite families are cyclic. Each satisfies \[
T_a\cap T_b=T_{\gcd(a,b)},\qquad
\gcd(|T_a|,|T_b|)=|T_{\gcd(a,b)}|.
\tag{1}
\] The binary \(r\) supplies the order-four complement in every odd
extension.

For any fixed odd base exponent \(e\), these two branches give precisely
the two possible base torus orders \(\gamma\). Choose the branch whose
order at \(e\) is \(\gamma\), and write its order at \(d\) as \(h_d\).
Since \(T_e\le T_{ea_i}\), \(\gamma\mid h_{ea_i}\). The other order at
\(ea_i\) is coprime to \(h_{ea_i}\). Because \(\gamma>1\), it cannot
also be divisible by \(\gamma\). Consequently each source condition
\(\gamma\ell_i\) dividing a nonsplit torus order automatically selects
this same compatible branch. The same argument identifies the prescribed
ambient branch. Thus all the source's choices can be made simultaneously
inside one \(\rho\)-stable torus.

For Construction 8 use the standard split tori \[
T_d=\{h(\lambda):\lambda\in\mathbb F_{2^d}^{\times}\},
\qquad h_d=2^d-1.
\] They satisfy the same intersection and gcd identities (1). Their
common complement is the binary \(w\), of order two, acting by
inversion. In the companion descent paper's standard matrix convention,
the diagonal subgroups are compatible under every odd subfield
embedding: when \(d\mid n\) and \(n/d\) is odd,
\((n+1)/2\equiv(d+1)/2\pmod d\), so the ambient exponent \(\theta_n\)
restricts to \(\theta_d\) on \(\mathbb F_{2^d}\).

\section{One intersection calculation covers both
constructions}\label{one-intersection-calculation-covers-both-constructions}

Write \(c_0=2\) for the split case and \(c_0=4\) for either nonsplit
case. Set \(\gamma=t-1\) in Construction 8; in Construction 9 retain the
chosen nonsplit base order. Thus \(\gamma=h_e\) in every case. In the
chosen ambient cyclic torus \(T_n\), let \[
A=C_m,\qquad m=\gamma\prod_i\ell_i,\qquad
K=A:C_{c_0}.
\] For Construction 9, \(m\mid h_n\) is a source assumption. For
Construction 8 it follows from the individual divisibilities: each
\(\gamma\ell_i\) divides \(2^n-1\), and, because the primes are
distinct, their least common multiple is \(m\). The common binary
complement normalizes \(A\).

Take the standard subfield parabolics
\(S_i=\operatorname{Sz}(2^{n/a_i})\). For \(J\subseteq\{1,\ldots,k\}\),
put \[
n_J=e\prod_{i\notin J}a_i,\qquad
m_J=\gamma\prod_{i\notin J}\ell_i.
\] Then \[
\bigcap_{j\in J}S_j=\operatorname{Sz}(2^{n_J}),\qquad
K\cap\bigcap_{j\in J}S_j=C_{m_J}:C_{c_0}.
\tag{2}
\]

Here is the valuation argument that handles \(\ell_i\mid\gamma\). For
the prime \(\ell_i\), the exponent in \(m\) is \(v_{\ell_i}(\gamma)+1\).
If \(i\notin J\), \(ea_i\mid n_J\), so \(\gamma\ell_i\mid h_{n_J}\); the
extra copy is retained. If \(i\in J\), then \[
\gcd(ea_i,n_J)=e,\qquad
\gcd(h_{ea_i},h_{n_J})=\gamma.
\] Since \(\gamma\ell_i\mid h_{ea_i}\), this forces
\(v_{\ell_i}(h_{n_J})=v_{\ell_i}(\gamma)\): exactly the extra copy is
removed. Primes not selected among the \(\ell_i\) occur in \(m\) only
through \(\gamma\), which is always retained. Thus
\(\gcd(m,h_{n_J})=m_J\). The binary complement belongs to every standard
subfield group, proving the full cyclic-by-complement equality in (2).

This proof allows \(a_i\mid e\): factoring \(e\) from both exponents
still gives \(\gcd(ea_i,n_J)=e\), since the remaining relative primes
are distinct from \(a_i\).

\section{Cyclic coset coordinates}\label{cyclic-coset-coordinates}

Let \(H_i=K\cap S_i=C_{m/\ell_i}:C_{c_0}\) and
\(A_i=A\cap H_i=C_{m/\ell_i}\). Every left coset of \(H_i\) in \(K\) has
a representative in \(A\), because the complement is contained in
\(H_i\), and \[
A/A_i\longrightarrow K/H_i,\qquad aA_i\longmapsto aH_i
\] is a bijection of sets.

The cyclic group \(A\) maps onto \(\prod_i A/A_i\): these are cyclic
quotients of the pairwise distinct prime orders \(\ell_i\), so this is
the ordinary Chinese remainder theorem. Its kernel has order \(\gamma\).
This remains true when some selected primes divide \(\gamma\). Hence the
set map \[
K\longrightarrow\prod_i K/H_i
\] is surjective. No \(H_i\) is asserted to be normal in \(K\). Only the
subgroups \(A_i\) of the abelian \(A\) are used as group-quotient
kernels.

The standard subfield types satisfy Theorem C of the companion descent
paper. Its Lemma D now gives a common representative in every flag of
the enlarged coset system. The system is a geometry, and \(G\) is
flag-transitive.

\section{RWPRI, thickness, and residual
connectedness}\label{rwpri-thickness-and-residual-connectedness}

The all-subset proof in the companion descent paper applies with (2).
Before selecting \(K\), each unused field type makes a maximal
simple-subfield step; their terminal group is
\(\operatorname{Sz}(2^e)\), with \(e\ge3\). Selecting \(K\) there gives
its full split or nonsplit torus normalizer, a maximal subgroup. After
selecting \(K\), an unused field type removes precisely one factor
\(\ell_i\), giving prime subgroup index. These are the source's RWPRI
conditions.

Panel sizes are the odd primes \(\ell_i\), and the index of the full
base torus normalizer in the simple group \(\operatorname{Sz}(2^e)\).
The latter is greater than two (an index-two subgroup of a simple group
would be normal). Thus the geometry is thick.

For residual connectedness, standard residues are enough because flag
transitivity is proved. With \(K\) already selected, two unused field
types give cyclic parts of orders \(m_J/\ell_i\) and \(m_J/\ell_j\).
Their least common multiple is \(m_J\), so they generate the current
cyclic-by-complement subgroup. Without \(K\), two remaining distinct
maximal simple subfields generate the current Suzuki group. If only one
field type and \(K\) remain, the simple base subfield is maximal, and
(2) shows that the other marked subgroup is not contained in it. They
therefore generate as well. This proves connectedness of every residue
of rank at least two, including the whole geometry.

We have thus rederived the complete parameterized \textbf{existence
assertions} of Constructions 8 and 9, and proved flag transitivity and
residual connectedness for these canonical realizations. Their inclusion
in an exhaustive classification, the classification of all embeddings
with those isomorphism types, and geometric intersection-property axioms
are separate questions.

\end{document}
